\documentclass[11pt]{article}
\usepackage{amsmath,amssymb,amsthm,bm}
\usepackage[margin=1in]{geometry}
\usepackage{booktabs}
\usepackage{array}
\usepackage{hyperref}
\usepackage{url}
\usepackage{listings}
\usepackage{xcolor}
\usepackage{longtable}

% Allow inline code identifiers to break at these characters so long
% file/function names never overrun their table column.
\makeatletter
\g@addto@macro\UrlBreaks{\do\_\do\:\do\-}
\makeatother

\definecolor{codebg}{RGB}{247,247,247}
\definecolor{codekw}{RGB}{0,0,160}
\definecolor{codecm}{RGB}{110,110,110}
\definecolor{codestr}{RGB}{160,0,0}

\lstset{
  basicstyle=\ttfamily\footnotesize,
  breaklines=true,
  frame=single,
  backgroundcolor=\color{codebg},
  columns=fullflexible,
  keepspaces=true,
  showstringspaces=false,
  keywordstyle=\color{codekw},
  commentstyle=\color{codecm}\itshape,
  stringstyle=\color{codestr},
  xleftmargin=0.5em,
  framexleftmargin=0.5em
}
% Inline code with literal underscores/slashes: \code|run_agas.py|
% Uses url's \path so long identifiers break instead of overrunning columns.
\urlstyle{tt}
\newcommand{\code}{\path}

\newtheorem{theorem}{Theorem}
\newtheorem{lemma}[theorem]{Lemma}
\newtheorem{proposition}[theorem]{Proposition}
\newtheorem{corollary}[theorem]{Corollary}
\theoremstyle{definition}
\newtheorem{definition}[theorem]{Definition}
\newtheorem{remark}[theorem]{Remark}
\newtheorem{example}[theorem]{Example}

\title{Attacking Recommender Systems with Agentic AI:\\
An Adaptive, Role-Switching Group Shilling Attack\\
\large A step-by-step walkthrough, from first principles to the code and the figures}
\author{}
\date{}

\begin{document}
\maketitle

\begin{abstract}
\noindent
This document explains the system in this code folder (\code|agent_attack_rs|)
from the very beginning. It assumes you know nothing. The code is the engine
behind the paper \emph{``An Efficient and Effective Agentic Group Shilling Attack
on Recommender Systems''} (AGAS). The compiled paper is \texttt{../main.pdf}.

The goal of the system is simple to state. We want to make one chosen item, the
\textbf{target item}, show up in the recommendation lists of ordinary users. We do
this by adding a small number of \textbf{fake users} to a recommender. What is new
is \emph{how} those fake users behave. A central \textbf{Coordinator}, which is a
language-model agent, watches how the platform reacts. Every round it does two
things. First it picks one \textbf{strategy}. Then it gives each fake user a
\textbf{role} out of four. The four roles are Profiler, Sniper, Camouflageur, and
Inactive. The roles change from round to round, and the strategy adapts when the
platform pushes back. Because of this the fake activity never settles into the
fixed pattern that detectors look for.

The document goes from the lowest level to the highest. We start with what a
recommender and a shilling attack even are (Part~I). We then look at the smallest
possible action, one fake rating (Part~II). We build up to one fake user with a
role (Part~III), then the signals the Coordinator reads (Part~IV), then the
Coordinator and its eight strategies (Part~V), then the victim models and how a
rank is measured (Part~VI), then the round loop that ties it all together
(Part~VII), and then the full command-line pipeline (Part~VIII). Part~IX shows how
the many output files turn into the single numbers printed in the paper. Part~X is
the reproduction guide. For every picture in \texttt{main.pdf} it says what the
picture shows, which command makes the numbers, and which script draws it. Part~XI
collects every equation and every file in one place.

Every symbol is defined before it is used. Every idea is tied to the exact file
and function that implements it. That way a new reader can trace ``paper to code''
without guessing.
\end{abstract}

\tableofcontents
\newpage


%======================================================================
\part*{Part I --- The problem, stated from nothing}
\addcontentsline{toc}{part}{Part I --- The problem, stated from nothing}

\section{What a recommender system is}

A recommender system watches what each user interacts with. It then tries to guess
what else that user would like. Write $\mathcal{U}=\{1,\dots,M\}$ for the users and
$\mathcal{I}=\{1,\dots,N\}$ for the items. The only thing we record is a big table
of zeros and ones,
\[
\mathbf{R}\in\{0,1\}^{M\times N},\qquad
R_{ui}=\begin{cases}1 & \text{user $u$ interacted with item $i$},\\ 0 & \text{otherwise.}\end{cases}
\]
This is called \textbf{implicit feedback}. We do not need star ratings. We only
need to know that an interaction happened. In this code the table is a pandas
\code|DataFrame| with three columns, \code|user_id|, \code|item_id|, and
\code|rating|. It is loaded and cleaned in
\code|src/agas/cli.py::_load_processed_tables|.

The recommender learns a scoring function $\hat{y}_{ui}=f_\Theta(u,i)$. This
function gives one real number for every pair of a user and an item. To make a
recommendation, it takes the items a user has not seen yet, sorts them by that
score, and returns the top ones:
\begin{equation}
\mathrm{TopK}(u)=\operatorname*{arg\,topK}_{q\in\mathcal{I}\setminus\mathcal{I}^{+}_u}\hat{y}_{uq},
\qquad
\mathcal{I}^{+}_u=\{i: R_{ui}=1\}.
\label{eq:topk}
\end{equation}
Here $\mathcal{I}^{+}_u$ is the set of items that user $u$ already interacted with.
This is Eq.~(topk\_def) in the paper's \texttt{formulation.tex}.

\begin{remark}[Which $f_\Theta$?]
The code has eleven concrete choices of $f_\Theta$. They include matrix
factorisation, neural collaborative filtering, and graph networks. They all live
under \code|src/agas/recsys/targets/|. There is also a fast lightweight surrogate
in \code|src/agas/recsys/surrogate.py| for the inner loop. We return to all of
these in Part~VI. For now, think of $f_\Theta$ as a black box that scores items.
\end{remark}

\section{What a shilling attack is}

A \textbf{shilling attack} is also called profile injection, or poisoning. The
idea is this. An attacker creates fake accounts and gives them fake interaction
histories. The recommender then retrains on this polluted data. After retraining,
one chosen \emph{target item} $i_{\mathrm{tar}}$ is recommended to more real users
than it deserves.

We can write this down. The attacker adds $M_f$ fake users. Each fake user $w$ has
its own interaction vector $\tilde{\mathbf{r}}_w\in\{0,1\}^N$. Stack these vectors
into a matrix $\tilde{\mathbf{R}}\in\{0,1\}^{M_f\times N}$. Then place that matrix
below the clean table:
\begin{equation}
\mathbf{R}'=[\mathbf{R};\,\tilde{\mathbf{R}}].
\label{eq:poison}
\end{equation}
This is Eq.~(poison\_append) in \texttt{formulation.tex}. In the code the
``placing below'' is just \code|pd.concat| followed by a refit. See
\code|src/agas/recsys/surrogate.py::append_interactions| and
\code|src/agas/recsys/targets/base.py::append_interactions|.

\section{How we measure success}

We measure success on a set of ordinary \emph{test} users
$\mathcal{U}_{\mathrm{test}}$. These are real users, and the attack never touches
them. The main question is a hit rate. Does the target item appear in a test
user's top-$K$ list?
\begin{equation}
\mathrm{HR}@K(i_{\mathrm{tar}})=\frac{1}{|\mathcal{U}_{\mathrm{test}}|}
\sum_{u\in\mathcal{U}_{\mathrm{test}}}\mathbb{I}\!\left[i_{\mathrm{tar}}\in\mathrm{TopK}(u)\right].
\label{eq:hr}
\end{equation}
There is also a rank-weighted version, NDCG@$K$. It gives more credit when the
item lands near the very top of the list:
\begin{equation}
\mathrm{NDCG}@K=\frac{1}{|\mathcal{U}_{\mathrm{test}}|}\sum_{u\in\mathcal{U}_{\mathrm{test}}}
\frac{\mathbb{I}[i_{\mathrm{tar}}\in\mathrm{TopK}(u)]}{\log_2(\operatorname{rank}(i_{\mathrm{tar}}\mid u)+1)}.
\label{eq:ndcg}
\end{equation}
These two are Eqs.~(hr\_k) and (ndcg\_k) in \texttt{experiment.tex}.

There is one more metric. It is used to check that the attack stays \emph{clean}.
It is Rec@$K$, the hit rate of the test users' \emph{own} held-out items:
\begin{equation}
\mathrm{Rec}@K=\frac{1}{|\mathcal{U}_{\mathrm{test}}|}\sum_{u\in\mathcal{U}_{\mathrm{test}}}
\mathbb{I}\!\left[i^{\mathrm{test}}_u\in\mathrm{TopK}(u)\right],
\label{eq:recK}
\end{equation}
where $i^{\mathrm{test}}_u$ is the item we hid from user $u$ for testing. This is
Eq.~(recK) in \texttt{experiment.tex}. The idea is easy. If Rec@$K$ barely drops
after the attack, then the recommender still serves real users well. That means
the attack did not visibly break the system.

\begin{remark}[A single-item shortcut used inside the code]
Equations~\eqref{eq:hr} and \eqref{eq:ndcg} average over many test users. That is
expensive to compute after every round. So inside the simulation the code uses a
cheaper stand-in. It tracks the \emph{mean rank of the target item} over a group
of receptive real users. It then turns one rank into HR and NDCG with a short
formula:
\begin{equation}
\mathrm{HR}@K=\mathbb{I}[\text{rank}\le K],\qquad
\mathrm{NDCG}@K=\frac{\mathbb{I}[\text{rank}\le K]}{\log_2(\text{rank}+1)}.
\label{eq:rank2hrndcg}
\end{equation}
This lives in \code|src/agas/cli.py::_rank_to_hr_ndcg|. Rank is the thing the
Coordinator can actually watch round by round, so it is the internal signal. The
HR and NDCG in the paper's big tables are the full averages of \eqref{eq:hr} and
\eqref{eq:ndcg}. We come back to this point in Part~IX because it matters when you
reproduce the numbers.
\end{remark}

\section{The threat model and the objective}

The attacker is \textbf{black-box}. It cannot see the victim's parameters, its
gradients, or its training code. It cannot see which users are test users, and it
cannot see their held-out items. All it sees is the ranked lists returned to
\emph{its own} fake accounts. This is stated in \texttt{formulation.tex}. In the
code it is enforced by \code|DefenseConfig(black_box_mode=True)| (Part~VI), which
hides the defender's private alerts from what the Coordinator observes.

Now put the goal and the stealth need together. The attacker wants to solve
\begin{equation}
\max_{\{\tilde{\mathbf{r}}_w\}}\;
\frac{1}{|\mathcal{U}_{\mathrm{test}}|}\sum_{u\in\mathcal{U}_{\mathrm{test}}}
\mathbb{I}\!\left[i_{\mathrm{tar}}\in\mathrm{TopK}(u;\mathcal{A}(\mathbf{R}'))\right]
\quad\text{s.t.}\quad \mathcal{S}(\tilde{\mathbf{R}};\mathbf{R})\le\varepsilon.
\label{eq:objective}
\end{equation}
Here $\mathcal{A}$ is the victim's training procedure. The term
$\mathcal{S}$ measures how strange the fake profiles look next to the real ones.
The tolerance $\varepsilon$ is never a fixed constant in the code. Instead it is
made real by the signals of Part~IV, which tell the Coordinator when to back off.
This is the whole point of the paper. Older attacks solve \eqref{eq:objective}
once, offline, with a fixed template. AGAS instead builds $\tilde{\mathbf{R}}$
step by step over $T$ rounds, and it adapts as the platform reacts.

\begin{center}
\fbox{\parbox{0.92\textwidth}{
\textbf{The whole system in one sentence.} A Coordinator agent adds fake ratings
round by round. It chooses a strategy and a role for each fake user from black-box
feedback. The target item's rank rises, and the fake profiles stay hard to tell
apart from real ones.}}
\end{center}


%======================================================================
\part*{Part II --- The lowest level: one fake rating}
\addcontentsline{toc}{part}{Part II --- The lowest level: one fake rating}

\section{The atom of the attack: a \texttt{RatingAction}}

Everything the attack does is built from one small thing. One fake user rates one
item. In the code this is a small dataclass,
\code|src/agas/agents/messages.py::RatingAction|. It carries \code|agent_id|,
which says which fake user acts. It carries \code|item_id|, which says which item.
It carries \code|rating|, which is the value. And it carries a free-text
\code|reason|, which is why the language model chose it.

\section{How a rating enters the world}

A proposed rating is not accepted right away. First it passes through the
environment's defense filter,
\code|src/agas/simulation/environment.py::_apply_action_with_defense|. This filter
models a realistic platform in three steps.
\begin{itemize}
\item It compares the rating to the item's current \textbf{bias} $b_i$, which is
just the item's average rating. The gap is $d=|r-b_i|$.
\item If the gap $d$ is large and the fake user has low trust, the rating is
\emph{discounted}. Its influence weight $w$ becomes less than one, and the
effective rating becomes $b_i + w\,(r-b_i)$.
\item If the platform is in ``lockdown'' and the user is not trusted, the rating
may be \emph{dropped} entirely, with some probability.
\end{itemize}
Only accepted, effective ratings are added to $\tilde{\mathbf{R}}$. The result of
each action, whether it was accepted, discounted, or dropped, is stored in an
\code|ActionOutcome|. These outcomes are the raw material for every signal in
Part~IV.

\begin{remark}[Why a rating can hurt on a graph victim]
On graph models such as LightGCN, every new edge raises the target's
\emph{degree}. The degree normalisation $D^{-1/2}AD^{-1/2}$ then weakens all of
the target's existing edges. So a naive ``rate the target 5.0'' can actually
\emph{lower} its rank. This one fact is the reason the attack uses \emph{bridge
items} (Part~VI) instead of rating the target directly on graph victims.
\end{remark}

\section{The poisoned matrix grows over time}

Over the rounds $t=0,\dots,T-1$, the accepted actions pile up:
\[
\tilde{\mathbf{R}}\leftarrow \tilde{\mathbf{R}}\cup\Delta\tilde{\mathbf{R}}^{(t+1)},
\qquad
\mathbf{R}^{(t)}=[\mathbf{R};\,\tilde{\mathbf{R}}^{(\le t)}].
\]
After each round the victim is refit on $\mathbf{R}^{(t)}$, and a fresh target
rank is read off. The final reported matrix is
$\mathbf{R}'=[\mathbf{R};\tilde{\mathbf{R}}]$ from \eqref{eq:poison}. The piling up
and the refit both happen in \code|environment.py::execute_step|, which calls
\code|append_interactions(..., refit=True)|.


%======================================================================
\part*{Part III --- One fake user and its four roles}
\addcontentsline{toc}{part}{Part III --- One fake user and its four roles}

\section{A worker is a fake user with a policy}

Each fake user is a \textbf{worker} agent,
\code|src/agas/agents/worker.py::WorkerAgent|. A worker holds a small mutable
state, \code|WorkerState|, which stores its trust, its risk, how many actions it
has taken, and its role history. Every round the Coordinator hands the worker a
\textbf{role}. The worker then turns that role into a list of
\code|RatingAction|s. It does this by prompting a language model in
\code|_act_with_llm|. If the model is not available, it falls back to simple
deterministic rules.

\section{The four roles \texorpdfstring{$\{$PR, SN, CA, IN$\}$}{PR SN CA IN}}

The paper uses exactly four roles. The enum is \code|src/agas/roles.py::Role|. For
the language-model protocol it is \code|messages.py::AgentRole|.

\begin{center}
\begin{tabular}{@{}llp{0.55\textwidth}@{}}
\toprule
Symbol & Name & What it does, and which item pool it may touch\\
\midrule
PR & Profiler & It makes a few safe filler ratings. It checks whether the platform
still accepts ratings, and it helps find \emph{bridge items}. Its pool is the
popular or benchmark items, the cluster items, or the bridge pool.\\
SN & Sniper & This is the payload role. On embedding victims it rates the target
directly. On graph victims it rates bridge or cluster items. It is given only to
workers with high trust and low risk. Its pool is the target plus competitors, or
the bridge pool.\\
CA & Camouflageur & This is the stealth role. It rates harmless, unrelated
``noise'' items so the fake user looks normal, and it rebuilds trust. Its pool is
the noise items.\\
IN & Inactive & It does nothing this round. It is used to cool a worker down or to
quarantine it.\\
\bottomrule
\end{tabular}
\end{center}

The map from a role to its allowed item pool and its action budget is in
\code|worker.py::_candidate_context|. The pools themselves are built by the
environment in \code|environment.py::build_worker_context|. These pools are the
benchmark, cluster, competitor, noise, bridge, flooding, and segment-profile
pools. After the language model answers, the host \emph{cleans} the reply in
\code|worker.py::_sanitize_llm_actions|. Any item outside the allowed pool is
thrown away, and every rating is clamped to the range $[1,5]$. This makes sure a
worker can only ever do what its role allows.

\section{How a role becomes a prompt}

Each role has a system prompt and a user prompt. They live under
\code|prompts/worker_*/| and are loaded by
\code|src/agas/llm/prompt_store.py::PromptStore|. If a file is missing, the code
falls back to defaults in \code|worker.py::_default_prompt_text|. Those defaults
are further tuned to the detected victim class, such as MF, LightGCN-small,
LightGCN-large, or Sequential. The worker fills a small JSON context. That context
holds its own signals, the allowed items, and a short memory of its own recent
outcomes. It sends this to the model through \code|src/agas/llm/providers.py| and
reads back a strict-JSON list of actions. The number of tokens used by every call
is added up for the efficiency figures.

\begin{center}
\fbox{\parbox{0.92\textwidth}{\textbf{Why role-switching gives stealth.} A fake
user that only ever snipes looks like a bot. A fake user that profiles, then
camouflages, then snipes once, then goes quiet, looks like a person. The
Coordinator changes roles every round, so no fake account ever settles into one
fixed signature. This is the mechanism behind the stealth results in Part~X.}}
\end{center}


%======================================================================
\part*{Part IV --- The signals the Coordinator reads}
\addcontentsline{toc}{part}{Part IV --- The signals the Coordinator reads}

Every decision the Coordinator makes is driven by two families of numbers. Both
live in \code|src/agas/signals.py| and match \code|method_coordinator.tex| in the
paper. This is the mathematical heart of the system, so we build each number from
scratch.

\section{Worker signals: trust \texorpdfstring{$\tau$}{tau}, risk \texorpdfstring{$\gamma$}{gamma}, validator \texorpdfstring{$\phi$}{phi}}

For each worker $w$ and round $t$ we keep three numbers. Let $r_{t,w}$ be an
accepted rating. Let $b_i$ be the item bias. Let $d_{t,w}=|r_{t,w}-b_i|$ be the
gap between them.

\paragraph{Trust $\tau_{t,w}$.} Trust goes up by one only when the accepted rating
looks benign. Benign means it stays within one point of the item's own average:
\begin{equation}
\tau_{t+1,w}=\max\!\bigl(0,\ \tau_{t,w}+\mathbb{I}[d_{t,w}\le 1]\bigr).
\label{eq:trust}
\end{equation}

\paragraph{Risk $\gamma_{t,w}$.} Risk goes up when the action is far from the bias.
It goes up even more when the action is extreme and the worker has not yet earned
trust:
\begin{equation}
\gamma_{t+1,w}=\max\!\Bigl(0,\ \gamma_{t,w}
+0.5\,\mathbb{I}[d_{t,w}>1]
+0.5\,\mathbb{I}[d_{t,w}>1.5\wedge\tau_{t,w}<1]\Bigr).
\label{eq:risk}
\end{equation}
If instead the worker stays inactive that round, its risk cools down:
\begin{equation}
\gamma_{t+1,w}=\max(0,\ \gamma_{t,w}-0.5).
\label{eq:riskdecay}
\end{equation}
An explicit alert or a block adds a further $+0.5$.

\paragraph{Validator $\phi_{t,w}$.} Trust and risk track \emph{recent actions}. The
validator instead scores the \emph{whole profile} for how suspicious it looks:
\begin{equation}
\phi_{t,w}=\min\!\bigl(1,\ 0.4\,e_{t,w}+0.3\,h_{t,w}+0.3\,o_{t,w}\bigr).
\label{eq:validator}
\end{equation}
Here $e$ is the fraction of extreme ratings. The flag $h\in\{0,1\}$ is one when the
worker gave a clear positive to the target. And $o$ is the fraction of the
worker's items that at least two other fake users also rated, which is a group
overlap. Equations \eqref{eq:trust} through \eqref{eq:validator} are implemented,
one method each, in \code|signals.py::WorkerSignals|. The three methods are
\code|update_after_action|, \code|apply_inactive_decay|, and
\code|update_validator|. The profile score also has a standalone version in
\code|src/agas/agents/profile_validator.py|.

\section{Environment signals: \texorpdfstring{$\rho,\Delta\rho,\eta,\xi,a$}{rho, delta-rho, eta, xi, a}}

These are round-level numbers. They describe how the \emph{platform} is reacting.

\paragraph{Rank feedback $\rho^{(t)}$ and $\Delta\rho^{(t)}$.} The number
$\rho^{(t)}$ is the mean target rank, and lower is better. The change is
$\Delta\rho^{(t)}=\rho^{(t-1)}-\rho^{(t)}$. A weak or negative $\Delta\rho$ tells
the Coordinator that the current plan has stalled.

\paragraph{Acceptance rate $\eta_t$.} This is the fraction of attempted actions
that were accepted this round. A sudden drop hints that the platform is quietly
filtering the attack.

\paragraph{Suppression signal $\xi_t=(q_t,s_t)$.} This is the composite the
Coordinator watches the most. Let $\mathcal{A}^{\mathrm{try}}_t,
\mathcal{A}^{\mathrm{drop}}_t,\mathcal{A}^{\mathrm{acc}}_t,\mathcal{A}^{\mathrm{tar}}_t$
be the attempted, dropped, accepted, and target-related actions of round $t$. From
these we build five normalised terms:
\begin{equation}
\begin{aligned}
\hat d_t &= \frac{|\mathcal{A}^{\mathrm{drop}}_t|}{\max(1,|\mathcal{A}^{\mathrm{try}}_t|)}
&&\text{(dropped ratio)},\\
\hat\delta_t &= \min\!\Bigl(1,\tfrac{1}{2|\mathcal{A}^{\mathrm{acc}}_t|}\textstyle\sum_{a}\delta(a)\Bigr)
&&\text{(discount magnitude)},\\
\hat m_t &= \frac{1}{\max(1,|\mathcal{A}^{\mathrm{tar}}_t|)}\textstyle\sum_a\mathbb{I}[\Delta\rho^{\mathrm{tar}}(a)<\epsilon_\rho]
&&\text{(weak target movement)},\\
\hat s_t &= \min(1,\ s_t/3) &&\text{(normalised streak)}.
\end{aligned}
\label{eq:suppterms}
\end{equation}
We also build a round-level group overlap. It is the largest pairwise Jaccard
overlap between the workers' rated-item sets:
\begin{equation}
g_t=\max_{u\ne v}\frac{|A_{t,u}\cap A_{t,v}|}{|A_{t,u}\cup A_{t,v}|}.
\label{eq:overlap}
\end{equation}
The suspicion score gives all five terms equal weight:
\begin{equation}
q_t=\min\!\bigl(1,\ 0.2\hat d_t+0.2\hat\delta_t+0.2\hat m_t+0.2\hat s_t+0.2g_t\bigr).
\label{eq:qscore}
\end{equation}
The streak $s_t$ goes up whenever the round shows filtering, discounting, weak
movement, or no rank gain. Otherwise it resets to zero.

\paragraph{Alert flag $a_t\in\{0,1\}$.} This is a coarse trigger. It is set to one
on an abnormal target spike, on very high overlap ($g_t\ge 0.8$), or on a sharp
drop in acceptance.

All of \eqref{eq:suppterms} through \eqref{eq:qscore}, together with $\eta_t$ and
$a_t$, are computed in one method, \code|signals.py::EnvSignals.update|. The
neutral start ($\eta_0=1,\ \xi_0=(0,0),\ a_0=0,\ \Delta\rho^{(0)}=0$) matches
line~5 of Algorithm~1 in the paper and \code|signals.py::initial_env_signals|.

\begin{remark}[Where the raw counts come from]
The attacker never sees the defender's private state. So the counts that feed
\eqref{eq:suppterms} are rebuilt from the black box. This is done by
\code|src/agas/agents/defender.py::DefenseMonitorAgent.analyze|. It turns hidden
accept, drop, and discount events into public ``signals'' that the Coordinator is
allowed to read. That split is what keeps the threat model of \eqref{eq:objective}
honest.
\end{remark}


%======================================================================
\part*{Part V --- The Coordinator and its eight strategies}
\addcontentsline{toc}{part}{Part V --- The Coordinator and its eight strategies}

\section{What the Coordinator does each round}

The Coordinator is the brain,
\code|src/agas/agents/coordinator.py::Coordinator|. Each round it does four things.
First it reads the observation, which is all the signals of Part~IV plus a short
memory. Second it selects one round-level \textbf{strategy}. Third it gives each
worker a role. Fourth, afterwards, it checks and gathers the accepted actions. The
language model does the first three inside \code|LLMCoordinatorPolicy.assign|.
After that, plain deterministic guardrails clean up the result, which we describe
below. This whole thing is Algorithm~1, the outer loop, and Algorithm~2, the
policy $\Pi$, in the paper.

\section{The eight strategies}

The enum is \code|src/agas/strategies.py::Strategy|. The plain descriptions are in
\code|STRATEGY_DESCRIPTIONS| and in \code|method_strategies.tex|.

\begin{center}
\begin{tabular}{@{}llp{0.62\textwidth}@{}}
\toprule
\# & Name & Trigger and behaviour\\
\midrule
S1 & Victim Probe & In the early rounds, profilers probe to decide whether the
victim is embedding-based or graph-based.\\
S2 & Bridge Building & For graph victims only. Build a pool of bridge items for the
snipers.\\
S3 & Warm-up & Build trust $\tau$ with one profiler and a few camouflageurs before
any sniper fires.\\
S4 & First Push & The first coordinated payload. A few snipers act while a
camouflageur stays on as cover.\\
S5 & Silent Slowdown & Triggered when $\eta_t$ drops but $a_t=0$. Reduce pressure,
and let high-trust workers act as camouflageurs.\\
S6 & Profile Cleanup & Triggered when $\phi$ is high, or $\xi$ is climbing, or
$a_t=1$. Freeze suspicious snipers and avoid overlap.\\
S7 & Safe Replacement & In a high-risk state, rotate risky workers out and bring in
rested, low-$\gamma$ ones.\\
S8 & Main Attack & In a calm state with the target still far, promote the best
worker to sniper and keep pushing.\\
\bottomrule
\end{tabular}
\end{center}

The language model emits a role for each worker. The code then \emph{infers} which
of S1 through S8 that pattern matches. This is done in
\code|coordinator.py::_infer_strategy|, so the choice can be counted later for the
ablation (Part~X).

\section{The victim-probe state machine and the guardrail overlays}

Two parts of the Coordinator are pure code, not the language model.

\paragraph{Probe and classify (S1).} Before the campaign,
\code|Coordinator._probe_classify| runs a small state machine. It rates the target
directly and watches $\Delta\rho$. If the direct signal helps, the victim is
embedding-based, which the code calls MF-style. If it hurts or does nothing, the
victim is graph-based, which the code calls LightGCN-style. An optional
``sequential'' probe tells apart recency-aware victims. The result sets
\code|victim_model_class|, which then steers every prompt and item pool.

\paragraph{Overlays (Algorithm~2, part ii).} After the language model proposes
roles, a set of guardrails is applied in \code|Coordinator.assign_roles|. They
stack on top of each other.
\begin{itemize}
\item \textbf{Suspicion lockout}, in \code|_apply_lockouts|. A sniper that recently
drew drops, discounts, or alerts is forced off the sniper role for a few steps.
\item \textbf{Validator guardrail}, in \code|_apply_validator_guardrail|. A worker
whose $\phi$ is above a threshold is forced INACTIVE.
\item \textbf{Profiler presence}, in \code|_ensure_profiler_presence|. This
guarantees a periodic scout probe even if the model forgot to add one.
\item \textbf{Co-occurrence bridging}, in \code|_apply_cooccurrence_bridging|. This
stamps the chosen bridge-item method into each assignment for graph victims.
\end{itemize}
These overlays are exactly the components the ablation study switches off (Part~X),
so each one has a stable name and a toggle.

\section{Complexity}

Let $T$ be the number of rounds. Let $|\mathcal{U}_f|$ be the number of workers.
Let $L$ be the per-worker budget, and let $V$ be the cost of one victim query. One
full campaign then costs
\[
O\!\bigl(T\,(c_C+|\mathcal{U}_f|(c_W+L)+V)\bigr).
\]
In words, there are $O(T)$ heavy Coordinator calls and $O(T|\mathcal{U}_f|)$ short
worker calls. This is the complexity section of \texttt{method.tex}. The key point
for the efficiency figures is that the cost grows linearly with the rounds and the
workers. There is no GAN, no bi-level solve, and no per-target surrogate to
retrain.


%======================================================================
\part*{Part VI --- The victims, and how a rank is measured}
\addcontentsline{toc}{part}{Part VI --- The victims, and how a rank is measured}

\section{The eleven victim models}

The victims live in \code|src/agas/recsys/targets/|. They share a base class,
\code|base.py::BaseTargetRecommender|. The base class handles the data indexing,
implicit-versus-explicit training, \code|rank_item|, and
\code|mean_scores_for_segment|. The concrete families are:
\begin{itemize}
\item \textbf{Matrix factorisation}, in \code|mf.py| (MF, GMF, BPR-MF, PMF, WMF,
NMF), \code|neumf.py| (NeuMF and NCF), and \code|svdpp.py|.
\item \textbf{Graph}, in \code|lightgcn.py| (LightGCN), \code|ngcf.py| (NGCF), and
\code|gcmc.py|.
\item \textbf{Others}, in \code|autorec.py| (autoencoders), \code|itemknn.py|, and
\code|sequential.py|.
\end{itemize}
The map from a name to a class is \code|cli.py::_build_target_model|. There is also
a fast CPU surrogate, \code|recsys/surrogate.py|, used for the inner loop when a
full torch victim would be too slow.

\section{Measuring the target rank}

We give the model a set of segment users and a set of candidate items. Then
\code|BaseTargetRecommender.rank_item| scores every candidate, sorts them, and
returns the 1-based position of the target. The candidate set is the target's
genre or keyword cluster, from
\code|environment.py::_resolve_target_cluster_items|. In transfer mode it can be
all items. This single number is the $\rho^{(t)}$ of Part~IV.

\section{Bridge items, and why graph victims need them}

We saw in Part~II that rating the target directly can hurt on a graph victim. So
the snipers rate \emph{bridge items} instead. A bridge item sits on a short graph
path from real target-fans to the target. Rating it lets graph diffusion carry the
signal to the target, without inflating the target's own degree. The selection
methods are in \code|environment.py::compute_bridge_items|. They include
co-occurrence, gradient, two-hop, low-degree and high-degree structural, and an
``auto'' method that chooses based on how many ratings the target already has. This
is strategy S2.


%======================================================================
\part*{Part VII --- The round loop (Algorithm 1)}
\addcontentsline{toc}{part}{Part VII --- The round loop (Algorithm 1)}

\section{The episode runner}

Everything above is tied together by
\code|src/agas/simulation/episode.py::AGASEpisodeRunner.run|. This function
\emph{is} Algorithm~1 of the paper. For each round $t$ it does the following.
\begin{enumerate}
\item It builds the observation $\mathbf{o}_t$ from the current signals and the
short memory, in \code|environment.observation|.
\item It optionally scores every worker's profile validator $\phi$.
\item It calls \code|coordinator.assign_roles|. This returns one strategy and a
role for each worker.
\item Each worker acts, in \code|worker.act|. This returns a batch of
\code|RatingAction|s.
\item It calls \code|environment.execute_step|. This applies the defense filter,
updates $\tau$ and $\gamma$, adds the accepted rows, refits the victim, and reads
the new rank.
\item It appends everything to \code|history|, and it stops early if the goal rank
is reached.
\end{enumerate}
At the very end it also computes victim-side embedding-cluster metrics, in
\code|src/agas/metrics/embedding_cluster.py|, for the stealth analysis. These
metrics are never shown to the attacker.

\section{Two ways to score the attack: Option A and Option B}

The transfer command, \code|cli.py::cmd_run_transfer|, supports two protocols.
\begin{itemize}
\item \textbf{Option A, offline transfer.} Run the whole campaign on an
\emph{episode} model, which is often the fast surrogate or a LightGCN. Then pull
out the accepted fake interactions, add them to the clean data, and train
\emph{fresh} victim models from scratch. Finally compare the target's rank on the
clean data against the poisoned data.
\item \textbf{Option B, in-loop.} Here the victim \emph{is} the episode model, and
it retrains every round. So the Coordinator gets true rank feedback. After the run,
a ``best-sequence retrain'' keeps only rounds $0$ through $t^\star$, where
$t^\star=\arg\min_t\rho^{(t)}$, and retrains once more. This is the $t^\star$
protocol noted in the README.
\end{itemize}
Both protocols write per-victim fields into the output JSON. Those fields are
\code|initial_rank|, \code|final_rank|, \code|best_rank|, \code|hr_at_k|, and
\code|ndcg_at_k|.


%======================================================================
\part*{Part VIII --- The end-to-end pipeline (how to run)}
\addcontentsline{toc}{part}{Part VIII --- The end-to-end pipeline (how to run)}

This part is the practical ``what do I type''. Every command is run from the repo
root, \code|agent_attack_rs/|.

\section{Install}

\begin{lstlisting}[language=bash]
pip install -e .            # core AGAS loop (CPU only)
pip install -e '.[targets]' # add torch victims (LightGCN, NeuMF, ...)
export OPENAI_API_KEY=sk-...   # required for the LLM Coordinator/workers
export OPENAI_MODEL=gpt-5.1    # optional, default model id
\end{lstlisting}
If \code|OPENAI_API_KEY| is not set, the language-model calls fail and the workers
fall back to deterministic behaviour. That is fine for a dry run, but the paper
numbers use the language model.

\section{Step 1 --- get and preprocess the data}

Drop each raw dataset into \texttt{data/<name>/}. The links are in
\texttt{README.md}. Then normalise everything into canonical
\code|interactions.csv| and \code|items.csv| files:
\begin{lstlisting}[language=bash]
python scripts/preprocess_all.py --data-root data --output-root processed
\end{lstlisting}
This calls \code|cli.py::cmd_preprocess|, which calls
\code|src/agas/data/pipeline.py|. The per-dataset loaders are in
\code|src/agas/data/loaders/|. The output lands in \texttt{processed/<name>/}. For
a quick smoke test, the tiny bundled \code|ml-latest-small| sample needs no
download.

\section{Step 2 --- run one attack episode}

The single entry point is \code|scripts/run_agas.py|. It is a thin wrapper over
\code|agas.cli run-episode|:
\begin{lstlisting}[language=bash]
python scripts/run_agas.py \
    --dataset ml-100k --victim lightgcn \
    --rounds 18 --seed 42 --n_workers 8 --budget 0.01 \
    --out outputs/agas_ml100k_lightgcn_seed42.json
\end{lstlisting}
The seven flags map onto the paper's protocol. Here $T$ is \code|--rounds|,
$|\mathcal{U}_f|$ is \code|--n_workers|, and $L$ is \code|--budget|. The run writes
a JSON trace, and next to it a human-readable \code|.log|. The log holds the
per-round ranks, the strategies, the roles, the token usage, and the role and
strategy activation tallies.

\section{Step 3 --- the four reviewer scripts, one per research question}

The \texttt{bash/} folder runs a tiny sweep for each research question. Each script
prints which figure or table it maps to.
\begin{lstlisting}[language=bash]
bash bash/run_performance.sh          # RQ1 -> tables/benchmark_unpopular.tex + radar
bash bash/run_stealth_and_detect.sh   # RQ2 + RQ3 -> t-SNE figs + detection table
bash bash/run_ablation.sh             # RQ4 -> ablation heatmaps + backbone table
bash bash/run_efficiency.sh           # RQ5 -> efficiency figures
\end{lstlisting}
For the full paper, widen the loops inside each script to all 6 datasets, 11
victims, and 5 seeds. The script \code|run_stealth_and_detect.sh| also runs the
detector battery through \code|agas.cli evaluate-detectors|.

\section{The output JSON, field by field}

A \code|run-episode| output carries several fields. It has \code|initial_rank|,
\code|final_rank|, and \code|best_rank| with its step. It has \code|history|, which
is the per-round observation, assignments, reports, and feedback. It has
\code|token_usage|, which is the coordinator, per-agent, and aggregate token
counts. It has \code|activation_stats|, which is the role and S1 through S8
strategy counts. And it has \code|embedding_cluster_metrics|. A
\code|run-transfer| output instead carries \code|option_a| and \code|option_b|.
Each of those is a dict keyed by victim with rank, HR, and NDCG before and after.
These are exactly the fields Part~IX reads.


%======================================================================
\part*{Part IX --- From many runs to one number}
\addcontentsline{toc}{part}{Part IX --- From many runs to one number}

\section{The reporting convention}

The paper reports the \textbf{mean and the 95\% confidence interval over 5 runs}.
The 5 runs are 5 random seeds. In this guide we use seeds \textbf{42, 43, 44, 45,
46}. Every promotion metric is also scaled by $10^3$. So a table cell that reads
``$40.0\pm0.2$'' means $\mathrm{HR}@10=0.0400\pm0.0002$. In other words, each cell
is the aggregate of 5 JSON files that differ only in the \code|--seed| flag.

\section{The aggregation helper}

The AGAS engine writes one JSON per run. On its own it does not produce the
aggregated cells. That last mile, from runs to numbers, is what
\code|scripts/aggregate_results.py| does. This script was added alongside this
document. It works in five steps.
\begin{enumerate}
\item It walks an output directory and loads every \code|*.json| inside it.
\item It recognises both schemas, the \code|run-episode| one and the
\code|run-transfer| one.
\item It turns each reported rank into HR@$K$ and NDCG@$K$ with the same formula
\eqref{eq:rank2hrndcg} the CLI uses, which is \code|_rank_to_hr_ndcg|.
\item It groups the runs by \code|(config, dataset, victim)| and prints the mean
and the 95\% CI, in the $\times 10^3$ convention.
\item It can also write a tidy CSV that a plotting script can read.
\end{enumerate}
\begin{lstlisting}[language=bash]
# aggregate a seed sweep into HR@10 / NDCG@10 / rank-delta with 95% CI
python scripts/aggregate_results.py outputs/multiseed_welltrained
# RQ1 sweep, also dump a CSV for the plot scripts
python scripts/aggregate_results.py outputs/rq1_performance --k 10 --csv summary.csv
# efficiency: also report mean tokens per run
python scripts/aggregate_results.py outputs/rq5_efficiency --tokens
\end{lstlisting}

\begin{remark}[An honest caveat: proxy versus population HR]
The released smoke outputs use the single-target-rank proxy
\eqref{eq:rank2hrndcg}. So their HR@10 is either 0 or 1 per run. It is not the
smooth population average \eqref{eq:hr} used in the big tables. For this reason the
aggregation script also reports \code|rank_delta|, which is how far the target
climbed. That number is stable and always defined. To reproduce the population
HR@$K$ exactly, you must evaluate \eqref{eq:hr} over the full
$\mathcal{U}_{\mathrm{test}}$ set on the retrained victim. The needed hooks,
\code|rank_item|, \code|mean_scores_for_segment|, and the segment users, are all
present in \code|base.py| and \code|environment.py|.
\end{remark}


%======================================================================
\part*{Part X --- Figure and table reproduction guide}
\addcontentsline{toc}{part}{Part X --- Figure and table reproduction guide}

This is the part to read when the question is simply ``which file do I run to get
picture $X$?''. Two facts frame everything.

\begin{center}
\fbox{\parbox{0.93\textwidth}{
\textbf{(1) The drawing scripts only draw.} They live in \texttt{../code/}, the
paper folder, and they are pure \texttt{matplotlib}. Given numbers, they render a
\texttt{.png} into \texttt{../figures/}. Several of them still keep their numbers in
a hard-coded array.\\[2pt]
\textbf{(2) The numbers come from this repo}, \code|agent_attack_rs|. You run the
engine (Part~VIII), then aggregate (Part~IX), then paste the aggregated values into
the drawing script's data block, or point the script at the CSV.}}
\end{center}

\noindent
The order to run things is always the same:
\[
\text{preprocess}\ \to\ \text{bash RQ script}\ \to\ \text{\texttt{aggregate\_results.py}}\ \to\ \text{drawing script in \texttt{../code/}}.
\]

\subsection*{Suggested first run}
Start with the smallest end-to-end path, just to see the whole loop work:
\begin{lstlisting}[language=bash]
python scripts/preprocess_all.py --data-root data --output-root processed
bash bash/run_performance.sh
python scripts/aggregate_results.py outputs/rq1_performance --csv summary.csv
python ../code/benchmark_popularity_regime_summary.py   # renders the radar
\end{lstlisting}

\subsection{Honest status: what this repo generates, and what it does not}
\label{sec:status}

Not every figure can be produced end to end by this repo alone, and it is only
fair to say so plainly. This repo implements the \textbf{AGAS} attack. It does
\emph{not} re-implement the ten baseline attacks (RandomAttack, Bandwagon, AUSH,
PoisonRec, PGA, GSPAttack, TargetedAttack, CLeaR, AgentSA, AgentAttack) or the
five detectors (BaseDetect, DHAGCN, PCASelectUsers, GAGE, MD-CBA). So the AGAS row
or column of every figure is generatable from the engine. The baseline and
detector entries are not, and this document does not fabricate them. They must come
from those methods' own code.

\begin{center}
\footnotesize
\renewcommand{\arraystretch}{1.2}
\setlength{\tabcolsep}{5pt}
\begin{tabular}{@{}p{0.34\textwidth}p{0.60\textwidth}@{}}
\toprule
Figure or table & Status in this repo\\
\midrule
Main table (\code|benchmark_unpopular.tex|) & AGAS rows: yes. Baseline rows: not implemented.\\
Popularity / intro radar & AGAS polygon: yes. Baseline polygons: not implemented.\\
Budget trade-off & Yes, fully (AGAS only figure).\\
Ablation strategy heatmap & Yes, fully (added the \code|--disable-strategies| flag).\\
Ablation component / size heatmap & Yes, fully (roles, memory, signals, validator, coordinator flags).\\
Efficiency rounds trade-off & HR@10 per round: yes. Time line: yes, now that runtime is recorded.\\
Efficiency method vs time & AGAS bar and line: yes. Baseline bars and lines: not implemented.\\
Tokens-until-Top-10 & AGAS tokens and seconds: yes. Baseline bars: not implemented.\\
Detector table (\code|detection_mf.tex|) & Not generatable: needs the 5 detectors and an \code|evaluate-detectors| command, neither present.\\
Benign Rec@10 bars & Not generatable: needs Rec@$K$, which needs a benign train/test split the pipeline does not create.\\
Stealth t-SNE (both panels) & Embeddings exist, but no export script yet. The drawing scripts currently simulate the points.\\
\bottomrule
\end{tabular}
\end{center}

\noindent
Three of these gaps are quick and I closed the first: (1) the efficiency
\emph{time} axes needed wall-clock timing, so the episode now records
\code|runtime_sec| (total) and \code|elapsed_sec|/\code|cumulative_sec| per round,
read with \code|aggregate_results.py --timing|. (2) Rec@$K$ and (3) the detector
battery are larger, because they need a benign split and the detector
implementations respectively, so they remain honest gaps rather than stubs. If you
want any specific baseline or detector implemented, that is a separate, scoped
piece of work.

\subsection{The one chain every quantitative figure shares}
\label{sec:chain}

Before we go figure by figure, here is the one chain that turns a single run into a
single plotted number. Every line, bar, and heatmap figure is just this chain,
evaluated at different settings and then arranged on axes. It has four steps.

\paragraph{Step 1: a rank is produced.} During a run, the victim is refit and the
target's position among the candidate items is read off. This is done by
\code|src/agas/recsys/targets/base.py::rank_item| (line~132). That function scores
every candidate, runs \code|torch.argsort(scores, descending=True)|, and returns
the 1-based position of the target $i_{\mathrm{tar}}$. Inside the loop it is called
through \code|environment.py::_target_rank| (line~277) and stored as
\code|self.current_rank| after every round, in \code|environment.py::execute_step|
(line~706). Call this integer $\rho$.

\paragraph{Step 2: the rank becomes HR and NDCG.} The CLI turns a rank into the two
promotion metrics with the formula \eqref{eq:rank2hrndcg}. This is
\code|cli.py::_rank_to_hr_ndcg| (line~124), where
$\mathrm{HR}@K=\mathbb{I}[\rho\le K]$ and
$\mathrm{NDCG}@K=\mathbb{I}[\rho\le K]/\log_2(\rho+1)$. In \code|cmd_run_transfer|
this happens in the \code|for k in metrics_ks| loop, at lines~1321--1326 for
Option~A and at lines~1462 and 1576 for the other paths. That loop writes
\code|hr_at_k| and \code|ndcg_at_k|, together with the clean-versus-poisoned
\code|initial_rank| and \code|final_rank| (lines~1307--1314), into the run's JSON.

\paragraph{Step 3: the runs are averaged.} We run each setting 5 times, once per
seed (42, 43, 44, 45, 46), so there are 5 JSON files. The script
\code|scripts/aggregate_results.py| reads those files. It re-applies the same
\code|rank_to_hr_ndcg| (line~44), groups the runs by
\code|(config, dataset, victim)| (line~198), and returns the mean and the 95\% CI
over the 5 seeds through \code|mean_ci| (line~51), scaled by $10^3$. Call the mean
$\bar y$ and the CI half-width $c$.

\paragraph{Step 4: the number is placed on an axis.} The mean $\bar y$ becomes a
bar height, or a line's $y$, or a heatmap cell value. The half-width $c$ becomes
the error bar. Which axis or cell it lands in depends only on the loop variable
that produced the run. That variable is the round $t$, the budget $b$, the ablated
component, the dataset, or the victim.

So for every figure below we only need three things. We need to know which loop
variable is swept. We need the command that runs it. And we need to know how the
resulting pair of a setting and its $\bar y$ is drawn. The efficiency figures
replace Step~2 with token and time counters. The tokens come from
\code|cli.py::_collect_token_usage| (line~615), in the field
\code|token_usage.aggregate.total_tokens|. The role and strategy tallies come from
\code|_compute_activation_stats| (line~661). The wall-clock comes from the run
timing.

\subsection{Performance: the main table and the radar (RQ1)}

\subsubsection*{Main table \texorpdfstring{\code|benchmark_unpopular.tex|}{benchmark unpopular} (Tab.~I)}
The loop variable swept is the triple (dataset, victim, seed). Run it like this:
\begin{lstlisting}[language=bash]
# widen the loops in run_performance.sh to all datasets/victims, 5 seeds
bash bash/run_performance.sh
python scripts/aggregate_results.py outputs/rq1_performance --k 10 --csv rq1.csv
\end{lstlisting}
Each printed line of \code|rq1.csv| is one cell block of the table. The column
\code|hr@10_x1000_mean| and its \code|hr@10_x1000_ci| form the ``$40.0\pm0.2$''
entry. The column \code|ndcg@10_x1000_mean| gives the NDCG entry. Steps 1 through 3
of the chain produce each number. The table just lays out the victim and the metric
in columns, and the method in rows. Here ``method'' means which config produced the
run. AGAS is the full system. The baselines are separate scripts, but their rows
are filled the same way, going from rank to HR and NDCG to the mean and CI.

\subsubsection*{Popularity radar \texorpdfstring{\code|benchmark_popularity_regime_summary.py|}{popularity radar}}
This figure shows three victim families, which are MF, graph-conv, and
graph-contrastive. In each family, every attack method is drawn as a filled
quadrilateral over four axes. The four axes are NDCG@10, HR@10, NDCG@50, and
HR@50.

The numbers come from RQ1 with head and mid targets across the three families.
After aggregation you get, for each method $m$ and family $f$, one vector
$v_{m,f}=[\text{NDCG@10},\text{HR@10},\text{NDCG@50},\text{HR@50}]$.

Now the drawing. In \code|benchmark_popularity_regime_summary.py| that vector is
the dict entry \code|RAW[m][f]|. For example, \code|RAW["AGAS"]["mf"]| is
\code|[58,70,70,78]|. The four axes are angles $\theta_k=2\pi k/4$ from
\code|radar_angles|, with \code|set_theta_offset(pi/2)| and direction $-1$. For one
method, vertex $k$ of its polygon is the polar point $(\theta_k,\ v_{m,f}[k])$. The
call \code|ax.plot(angles, values)| draws the outline, and \code|ax.fill| shades
it, both inside \code|draw_radar|. So on the radar the angle tells you which metric
you are on, and the radius tells you the value of that metric. To make the figure
data-driven, replace the \code|RAW| block with your aggregate. The intro radar,
\code|intro_fig_radar_motivation.py|, is built the same way, but it is split by
popularity regime instead of by victim family.

\subsection{Stealthiness: the two t-SNE panels (RQ2)}

\subsubsection*{Left panel \texorpdfstring{(\code|benchmark_fig7_tsne_distance_users_target.py|)}{}}
This panel has one small plot per method, six in all. It samples nine users. Next
to each user it prints two numbers. The red number is the distance from that user
to the target item. The black number is the user's average distance to its own
ground-truth items. A method is best when both numbers are small. That means the
target was pulled close to the user, and the user was not dragged away from its
real tastes.

These two numbers come from a poisoned victim's learned vectors. After a run, the
surrogate exposes \code|user_factors| and \code|item_factors| in
\code|recsys/surrogate.py|. Embed the target item $i_{\mathrm{tar}}$, each sampled
user $u$, and the items $u$ rated. Call the red number $A_u$ and the black number
$B_u$. Then $A_u=\lVert e_u-e_{i_{\mathrm{tar}}}\rVert$ and
$B_u=\frac{1}{|\mathcal I^{+}_u|}\sum_{i\in\mathcal I^{+}_u}\lVert e_u-e_i\rVert$.
The way to reach the victim's vectors mirrors
\code|src/agas/metrics/embedding_cluster.py|, which already pulls the factor
matrices. That file uses cosine distance, and here we use Euclidean, but the hook
is the same.

Now an important caveat about the drawing. In this script the \emph{positions} of
the stars are not the distances. They are a stylised layout, set by
\code|ANGLES_DEG| and the per-method \code|LAYOUT[radius]|. The layout is the same
across panels so the eye can compare spread. The real data are only the printed
numbers. They live in the dict entry
\code|PROFILES[method][user]|, which is the pair $(A_u,B_u)$. They are drawn by
\code|ax.text(...)| at the star for the black number and at the connector midpoint
for the red number, inside \code|plot_panel|. So to make it real, replace the
\code|PROFILES| block with your measured $(A_u,B_u)$. The layout stays decorative
on purpose.

\subsubsection*{Right panel \texorpdfstring{(\code|dappoison_fig6_tsne.py|)}{}}
This panel has one 2-D t-SNE scatter per method. Blue dots are clean users, and red
stars are injected fake users. Diffuse red is stealthy, and clustered red is easy
to detect.

The numbers come from a poisoned matrix $\mathbf R'$. Get each user's history
embedding from the victim, then run \code|sklearn.manifold.TSNE| down to two
dimensions. The clean points are the real users. The fake points are the rows of
$\tilde{\mathbf R}$. The fake user ids are \code|agent_1..agent_N|, and the
injected rows come from \code|cli.py::_extract_attack_interactions|. Each point is
one $(x,y)$, its t-SNE coordinate. The call \code|ax.scatter| plots blue circles
for the clean users and red stars for the fake ones, inside \code|plot_panel|. One
note on the current state. The script simulates the two clouds with Gaussian
generators such as \code|fakes_agas| and \code|fakes_poisonrec|, so it renders
without a run. To make it data-driven, feed the real 2-D coordinates in place of
those generators.

\subsection{Detector: the table and the Rec@10 bars (RQ3)}

\subsubsection*{Detection table \texorpdfstring{\code|detection_mf.tex|}{detection}}
The loop variable swept is the pair (dataset, method). Run
\code|bash bash/run_stealth_and_detect.sh|. It first produces a poisoned trace,
then scores that trace with
\code|agas.cli evaluate-detectors --detectors BaseDetect,DHAGCN,PCASelectUsers,GAGE,MD-CBA|.
Each detector returns four numbers, which are Accuracy, Recall, Precision, and F1,
on the fake-versus-real labels. Those four numbers become the four sub-columns of
one block, for each pair of a detector and a method. Lower numbers mean the attack
is harder to detect.

\subsubsection*{Benign Rec@10 bars \texorpdfstring{(\code|detect_fig_defense_k5.py|)}{}}
This figure has three victims, which are NeuMF, LightGCN, and LightCCF. Each method
gets one horizontal bar. The solid part of the bar is the Rec@10 after the attack.
The hatched remainder is the no-attack ceiling. The printed percentage is the drop
in utility.

There are two numbers per bar. Call them $A$ and $B$. The value $B$ is the Rec@$K$
of the clean victim, the one with no injected users. The value $A$ is the Rec@$K$
of the poisoned victim. Rec@$K$ is \eqref{eq:recK}. To get it, evaluate each test
user's held-out item on the retrained model, using the same \code|rank_item|
machinery but scoring the user's own item instead of the target. So $B$ uses the
clean fit and $A$ uses the poisoned fit. These are exactly the clean and attacked
pair that the transfer command already builds, namely \code|clean_model| and
\code|attacked_model| in \code|cli.py| lines~1305--1316.

Now the drawing. In the script, \code|NO_ATTACK[victim]| is $B$ and
\code|AFTER_ATTACK[victim][method]| is $A$. Take the method at row $y_i$. The call
\code|ax.barh(y, before)| draws the full hatched bar, whose length is $B$. Then
\code|ax.barh(y, after)| draws the solid bar on top, whose length is $A$. The
leftover gap $B-A$ is the grey segment. Finally
\code|ax.text(..., f"{100*(B-A)/B:.0f}%")| prints the drop, all inside
\code|draw_panel|. So the $x$ position is the Rec@10 value, the $y$ position is the
method index, the solid bar length is $A$, and the ceiling is $B$.

\subsection{Ablation: two heatmaps and a table (RQ4)}

First, a word on notation. In these figures ``w/o'' just means ``without''. Each
``w/o~X'' row is one run of AGAS with component X switched off. We compare that run
to the ``Full'' run, which has everything on. If the number drops a lot, then X
mattered.

\subsubsection*{The full command for every row}
Every ablation row is a real flag on \code|scripts/run_agas.py|. Nothing here needs
a source edit. This one block runs the entire RQ4 sweep, all rows, across the 5
seeds. Copy, paste, run.
\begin{lstlisting}[language=bash]
DATASET=ml-100k ; VICTIM=lightgcn
COMMON="--dataset $DATASET --victim $VICTIM --rounds 18 --n_workers 8 --budget 0.01"
mkdir -p outputs/rq4

for seed in 42 43 44 45 46; do
  # Full system (everything on) -- the baseline the "w/o" rows are compared to
  python scripts/run_agas.py $COMMON --seed $seed \
      --coordinator-agent-memory --profile-validator \
      --out outputs/rq4/full__seed${seed}.json

  # ----- component / size heatmap rows -----
  python scripts/run_agas.py $COMMON --seed $seed --random-coordinator      --out outputs/rq4/wo_coordinator__seed${seed}.json
  python scripts/run_agas.py $COMMON --seed $seed --disable-roles pr        --out outputs/rq4/wo_profiler__seed${seed}.json
  python scripts/run_agas.py $COMMON --seed $seed --disable-roles sn        --out outputs/rq4/wo_sniper__seed${seed}.json
  python scripts/run_agas.py $COMMON --seed $seed --disable-roles ca        --out outputs/rq4/wo_camouflageur__seed${seed}.json
  python scripts/run_agas.py $COMMON --seed $seed --disable-roles in        --out outputs/rq4/wo_inactive__seed${seed}.json
  python scripts/run_agas.py $COMMON --seed $seed --no-coordinator-agent-memory --out outputs/rq4/wo_memory__seed${seed}.json
  python scripts/run_agas.py $COMMON --seed $seed --disable-signals         --out outputs/rq4/wo_signals__seed${seed}.json
  python scripts/run_agas.py $COMMON --seed $seed --no-profile-validator    --out outputs/rq4/wo_validator__seed${seed}.json

  # ----- strategy heatmap rows -----
  for s in s1 s2 s3 s4 s5 s6 s7 s8; do
    python scripts/run_agas.py $COMMON --seed $seed --disable-strategies $s \
        --out outputs/rq4/wo_${s}__seed${seed}.json
  done
done

# turn the 5-seed runs into one mean +/- 95% CI per row
python scripts/aggregate_results.py outputs/rq4 --k 10 --csv rq4.csv
\end{lstlisting}
The wrapper \code|bash bash/run_ablation.sh| runs exactly this (set
\code|SEEDS="42"| for a quick one-seed smoke). To read one row on its own, just take
its line. For example, the exact command for ``w/o~S7'' is
\begin{lstlisting}[language=bash]
python scripts/run_agas.py --dataset ml-100k --victim lightgcn \
    --rounds 18 --n_workers 8 --budget 0.01 --seed 42 \
    --disable-strategies s7 --out outputs/rq4/wo_s7__seed42.json
\end{lstlisting}
and the exact command for ``w/o~Profiler'' is the same line with
\code|--disable-roles pr| in place of \code|--disable-strategies s7|.

\subsubsection*{What each row removes (plain words)}
\begin{center}
\footnotesize
\renewcommand{\arraystretch}{1.2}
\setlength{\tabcolsep}{5pt}
\begin{tabular}{@{}llp{0.46\textwidth}@{}}
\toprule
Row & Flag & What it takes away\\
\midrule
w/o Coordinator & \code|--random-coordinator| & roles are chosen at random, not by the LLM\\
w/o Profiler & \code|--disable-roles pr| & no probing or bridge discovery (those workers become camouflageurs)\\
w/o Sniper & \code|--disable-roles sn| & no targeted payload (snipers become camouflageurs)\\
w/o Camouflageur & \code|--disable-roles ca| & no benign filler cover (camouflageurs go inactive)\\
w/o Inactive & \code|--disable-roles in| & no cool-down, everyone always acts\\
w/o Memory & \code|--no-coordinator-agent-memory| & the Coordinator forgets earlier rounds\\
w/o Signals & \code|--disable-signals| & trust, risk, validator, suppression, and alerts are hidden (only raw rank remains)\\
w/o Validator & \code|--no-profile-validator| & the profile-suspicion guardrail is off\\
w/o S1 & \code|--disable-strategies s1| & no Victim Probe (skip the probe phase)\\
w/o S2 & \code|--disable-strategies s2| & no Bridge Building\\
w/o S3 & \code|--disable-strategies s3| & no Warm-up (prompt-level)\\
w/o S4 & \code|--disable-strategies s4| & no First Push (prompt-level)\\
w/o S5 & \code|--disable-strategies s5| & no Silent Slowdown (prompt-level)\\
w/o S6 & \code|--disable-strategies s6| & no Profile Cleanup guardrail\\
w/o S7 & \code|--disable-strategies s7| & no Safe Replacement (no suspicion lockout)\\
w/o S8 & \code|--disable-strategies s8| & no Main Attack push (prompt-level)\\
\bottomrule
\end{tabular}
\end{center}
Two honest notes. First, you can combine flags, for example
\code|--disable-roles pr,sn| or \code|--disable-strategies s1,s7|. Second, S1, S2,
S6, S7 and every component row are hard switches that always take effect. S3, S4,
S5, S8 are behaviours the language model produces on its own, so ``w/o'' for those
is done by telling the Coordinator in its prompt not to use them. The model usually
obeys but is not forced.

\subsubsection*{Strategy heatmap \texorpdfstring{(\code|ablation_strategy_heatmap.py|)}{}}
This figure has four panels, one per dataset. In each panel the rows are Full, then
w/o~S1 through w/o~S8. The columns are three victims, which are NeuMF, LightGCN,
and LightCCF. The colour of a cell, and the number in it, is HR@10. You produce the
rows with the \code|--disable-strategies s1 ... s8| runs above.

Now the drawing. The script holds a matrix \code|VALUES[dataset]| of shape
$9\times3$. Row $i$ is the ablation config, and column $j$ is the victim. The call
\code|ax.imshow(values, cmap="YlOrRd_r")| colours cell $(i,j)$ by its HR@10, and
\code|ax.text(j, i, f"{val:.1f}")| prints the number, both inside
\code|draw_panel|. So the cell at row $i$, column $j$ is the HR@10, scaled by
$10^3$. Darker or brighter follows the colour bar. To fill it with your own runs,
replace the \code|VALUES| block with the aggregate from
\code|aggregate_results.py|.

\subsubsection*{Component and size heatmap \texorpdfstring{(\code|ablation_size_heatmap.py|)}{}}
The drawing is identical to the strategy heatmap. What changes is the axes. The rows
are now Full, w/o~Coordinator, w/o~Profiler, w/o~Sniper, w/o~CA, w/o~IN,
w/o~Memory, and w/o~Signals, which are the component runs above. The columns are the
backbone size, which is Large, Medium, or Small. You set the size by adding
\code|--worker-llm-model| to any of those commands. For example, one Large cell for
the ``w/o~Sniper'' row is:
\begin{lstlisting}[language=bash]
python scripts/run_agas.py --dataset ml-100k --victim lightgcn \
    --rounds 18 --n_workers 8 --budget 0.01 --seed 42 \
    --disable-roles sn --worker-llm-model llama-3.3-70b \
    --out outputs/rq4/wo_sniper__large__seed42.json
\end{lstlisting}
Use a 14B model for Medium and a 2B model for Small. So one cell is: run AGAS with
this size and this component removed, then read the aggregated HR@10.

\subsubsection*{Backbone table \texorpdfstring{\code|ablation_backbones.tex|}{backbones}}
This uses the same size sweep over language-model backbones, grouped as frontier,
large, medium, and small. You set the backbone with \code|--worker-llm-model| (and
\code|--llm-model| for the Coordinator). Each cell is HR@10, NDCG@10, and Rec@50,
aggregated per pair of a backbone and a method.

\subsection{Efficiency: rounds, time, tokens, and budget (RQ5)}

\subsubsection*{Rounds trade-off \texorpdfstring{(\code|efficiency_rounds_tradeoff.py|)}{}}
This is a plot with two y-axes. The bars are HR@10 per round, on the left axis. The
red line is the cumulative minutes per round, on the right axis. A vertical dashed
line marks the chosen budget, $T=18$.

The loop variable swept is the round index $t$. Inside one run, the per-round rank
is already in \code|history[t].feedback.target_rank|. Convert each rank with
\code|_rank_to_hr_ndcg| to get HR@10 at round $t$, then average over seeds for the
bar height and its CI. The line value is the wall-clock elapsed up to round $t$,
which the episode now records as \code|cumulative_sec| in each \code|history| entry
(and \code|runtime_sec| for the whole run). Read the totals with
\code|aggregate_results.py --timing|.

Now the drawing. The arrays \code|HR10[t]|, \code|HR10_CI[t]|, and
\code|TIME_MIN[t]| are indexed by \code|ROUNDS|, which runs from 1 to 30. The call
\code|ax1.bar(ROUNDS, HR10)| places a bar at $x=t$ of height \code|HR10[t]|, and
\code|ax1.errorbar(...)| adds whiskers of size \code|HR10_CI[t]|. The twin axis
call \code|ax2.plot(ROUNDS, TIME_MIN)| draws the point $(t,\text{minutes}(t))$.
Finally \code|ax1.axvline(18)| marks the budget. So a bar is the pair
$(x=t,\ y=\overline{\mathrm{HR}@10}(t))$ and a line point is the pair
$(x=t,\ y=\overline{\text{minutes}}(t))$.

\subsubsection*{Method versus time \texorpdfstring{(\code|efficiency_method_time_compare.py|)}{}}
This shows four datasets as x-groups, with five methods in each group. Each method
gets a bar, which is HR@10 on the left axis, and a line marker, which is the total
runtime in minutes on the right axis.

The loop variable swept is the pair (method, dataset). The bar height is the
aggregated HR@10 for that pair, from Steps 1 through 3. The line value is the mean
total minutes of those runs.

Now the drawing. The arrays \code|HR10[m]| and \code|TIME_MIN[m]| are length-4
lists over the datasets. With \code|x=arange(4)| and a per-method \code|offsets|,
the call \code|ax1.bar(x+offset[m], HR10[m])| puts method $m$'s bar for dataset $d$
at $x=d+\text{offset}(m)$, with height \code|HR10[m][d]|. The twin call
\code|ax2.plot(x+offset[m], TIME_MIN[m])| draws its runtime marker at the same $x$.
So the bar height is the effectiveness, and the marker height is the cost, read on
opposite axes.

\subsubsection*{Tokens until Top-10 \texorpdfstring{(\code|efficiency_token_topk.py|)}{}}
This shows grouped bars on a log y-axis. The bar is the average number of
language-model tokens spent before the target first enters the Top-10. There are
three agentic methods across four datasets. A side table lists the inference
seconds.

Here is how to get the numbers. Walk a run's \code|history| to the first round
$t^\ast$ where \code|feedback.target_rank| is at most 10. Then sum the per-call
tokens up to $t^\ast$. Each report carries \code|trace.token_usage.total_tokens|,
and the run total is \code|token_usage.aggregate.total_tokens| from
\code|_collect_token_usage| (line~615). Average this over seeds. The seconds column
is the matched wall-clock. The command
\code|python scripts/aggregate_results.py <dir> --tokens| prints the per-run token
mean that feeds this figure.

Now the drawing. The arrays \code|TOKENS[m]| and \code|TIME_S[m]| are length-4
lists. The call \code|ax_bar.bar(x+offset[m], TOKENS[m])| with
\code|ax_bar.set_yscale("log")| makes a bar at $x=d+\text{offset}(m)$ whose height,
on a log scale, is the token count. The call \code|ax_bar.text(...)| prints a label
such as ``55K'' above it. The right-hand \code|ax_tab.table| lists \code|TIME_S|
row by row, dataset first and method within. So the bar height, on a log scale, is
the tokens, and the table cell is the seconds.

\subsubsection*{Budget trade-off \texorpdfstring{(\code|efficiency_fig_agents_tradeoff.py|)}{}}
This has a 2-by-2 grid of subplots, one per dataset. In each subplot there are four
lines, which are HR@10, NDCG@10, HR@50, and NDCG@50. The x-axis is the fake-user
budget, from 0.5\% to 3.0\%.

The loop variable swept is the budget $b$. The script \code|run_efficiency.sh|
loops \code|--budget 0.005 0.01 0.02 0.03|, and you can add 0.015 and 0.025. Each
budget value is one x-tick. Aggregate HR and NDCG per budget.

Now the drawing. The entry \code|VALUES[dataset][metric]| is a length-6 list over
\code|BUDGETS=[0.5,1,1.5,2,2.5,3]|. For each metric, the call
\code|ax.plot(BUDGETS, VALUES[dataset][metric])| draws the point
$(x=b,\ y=\text{metric}(b))$ and connects the points, one styled line per metric.
So each line is one metric as the budget grows, and the four lines share the
subplot.

\subsection{Tables at a glance}

\begin{center}
\footnotesize
\renewcommand{\arraystretch}{1.25}
\setlength{\tabcolsep}{4pt}
\begin{tabular}{@{}p{0.20\textwidth}p{0.30\textwidth}p{0.42\textwidth}@{}}
\toprule
Table (in \texttt{main.pdf}) & Produced by (this repo) & Numbers from\\
\midrule
\code|benchmark_unpopular.tex| (Tab.~I, main result) &
\code|bash/run_performance.sh| over all datasets and victims &
\code|option_a| or \code|option_b| ranks turned into HR and NDCG by
\code|aggregate_results.py|\\
\code|benchmark_popularity_regime.tex| &
the same, with head and mid target items &
the same aggregation on head and mid target bins\\
\code|detection_mf.tex| (detector table) &
\code|bash/run_stealth_and_detect.sh|, then \code|evaluate-detectors| &
detector Accuracy, Recall, Precision, and F1 per method\\
\code|ablation_mf.tex|, \code|ablation_backbones.tex| &
\code|bash/run_ablation.sh| with component toggles, plus \code|--worker-llm-model|
for the backbones &
HR and NDCG per ablated config, aggregated over seeds\\
\code|efficiency_breakdown.tex| &
\code|bash/run_efficiency.sh| &
per-stage wall-clock from the run timing\\
\bottomrule
\end{tabular}
\end{center}

\begin{remark}[Illustrative versus data-driven, said plainly]
Today several drawing scripts carry a hard-coded data block. These are \code|RAW|,
\code|VALUES|, \code|HR10|, \code|TOKENS|, \code|PROFILES|, and the t-SNE
generators. That block is only the last inch of Step~4. It has exactly the shape
that the aggregation in Part~IX prints. So making a figure data-driven means one
thing. You paste the aggregated array, or read the \code|--csv|, in place of the
hard-coded one. Nothing about the geometry changes. The engine in this repo owns
every number, and the scripts in \texttt{../code/} own only the drawing.
\end{remark}


%======================================================================
\part*{Part XI --- Everything in one place}
\addcontentsline{toc}{part}{Part XI --- Everything in one place}

\section{Every equation, and where it lives}

\begin{center}
\footnotesize
\setlength{\tabcolsep}{4pt}
\begin{tabular}{@{}llp{0.46\textwidth}@{}}
\toprule
Quantity & Equation & File and symbol\\
\midrule
Top-$K$ recommendation & \eqref{eq:topk} & \code|recsys/*::recommend| and \code|rank_item|\\
Poisoned matrix & \eqref{eq:poison} & \code|*::append_interactions|\\
HR@$K$, NDCG@$K$ & \eqref{eq:hr},\eqref{eq:ndcg} & \code|experiment.tex| (population)\\
rank to HR and NDCG & \eqref{eq:rank2hrndcg} & \code|cli.py::_rank_to_hr_ndcg|\\
Attack objective & \eqref{eq:objective} & \code|formulation.tex|\\
Trust $\tau$ & \eqref{eq:trust} & \code|signals.py::update_after_action|\\
Risk $\gamma$ & \eqref{eq:risk},\eqref{eq:riskdecay} & \code|signals.py::update_after_action| and \code|apply_inactive_decay|\\
Validator $\phi$ & \eqref{eq:validator} & \code|signals.py::update_validator|\\
Suppression terms & \eqref{eq:suppterms} & \code|signals.py::EnvSignals.update|\\
Group overlap $g_t$ & \eqref{eq:overlap} & \code|signals.py::EnvSignals.update|\\
Suspicion $q_t$ & \eqref{eq:qscore} & \code|signals.py::EnvSignals.update|\\
\bottomrule
\end{tabular}
\end{center}

\section{Every file that matters}

\begin{center}
\footnotesize
\setlength{\tabcolsep}{4pt}
\begin{tabular}{@{}p{0.40\textwidth}p{0.54\textwidth}@{}}
\toprule
File & Role\\
\midrule
\code|scripts/run_agas.py| & single entry point for one episode\\
\code|scripts/preprocess_all.py| & dataset into canonical CSVs\\
\code|scripts/aggregate_results.py| & runs into aggregated numbers (added by this doc)\\
\code|bash/run_*.sh| & one sweep per research question\\
\code|src/agas/roles.py|, \code|strategies.py| & the 4 roles and the 8 strategies\\
\code|src/agas/signals.py| & all worker and environment signals\\
\code|src/agas/agents/coordinator.py| & strategy choice, role assignment, guardrails\\
\code|src/agas/agents/worker.py| & per-role action generation\\
\code|src/agas/agents/defender.py| & black-box signal reconstruction\\
\code|src/agas/simulation/episode.py| & the round loop (Algorithm 1)\\
\code|src/agas/simulation/environment.py| & defense filter, refit, rank, item pools, bridges\\
\code|src/agas/recsys/targets/*| & the 11 victim models\\
\code|src/agas/metrics/embedding_cluster.py| & victim-side stealth metric\\
\code|prompts/*| & Coordinator and per-role prompt templates\\
\bottomrule
\end{tabular}
\end{center}

\section{The whole pipeline in one paragraph}

First preprocess the datasets into canonical CSVs. Then run
\code|scripts/run_agas.py|, or a \code|bash/run_*.sh| sweep, to play out $T$
rounds. In each round a Coordinator language model reads the trust, risk, and
validator signals, and the rank, acceptance, suppression, and alert signals. It
picks one of eight strategies and gives each fake user one of four roles.
Deterministic guardrails keep the campaign stealthy. Each round the victim is
refit and the target's rank is read again. This produces one JSON trace per run.
Next, aggregate those traces across five seeds with
\code|scripts/aggregate_results.py| into HR@$K$ and NDCG@$K$, reported as the mean
and 95\% CI and scaled by $10^3$. Finally feed those numbers to the matplotlib
scripts in \texttt{../code/} to draw every figure in \texttt{main.pdf}.

\end{document}
